\documentclass[11pt,a4paper]{article}
\usepackage[margin=2.4cm]{geometry}
\usepackage[T1]{fontenc}
\usepackage[hidelinks]{hyperref}
\setlength{\parskip}{5pt}\setlength{\parindent}{0pt}
\begin{document}
\begin{flushright}Danang, Vietnam\\ \today\end{flushright}

Editor-in-Chief\\
Computer Communications

\textbf{Submission of ``Training Budget, Time Parameterisation and Link Pattern Decide the Ranking of
Propagation Operators on LEO Inter-Satellite-Link Graphs''}

Dear Editor,

I submit the manuscript above for consideration as a regular paper in Computer Communications.

\textbf{What the paper is.} It is an audit, with a pre-declared protocol, of a result I reported
myself in a preliminary version circulated for review in June 2026: that the transition-probability
kernel of a continuous-time quantum walk outperforms the heat kernel as a propagation operator on a
1584-satellite low-Earth-orbit inter-satellite-link graph. Three operators sharing one eigenbasis and
one network are compared on Walker-delta shells of 264 to 4400 satellites, two node-field tasks and
ten seeds, across arms that vary the parameterisation and initialisation of the learned propagation
time, the capacity, the training budget and the link pattern, against fixed-reach baselines whose
reach is measured: 5590 training runs in all. The preliminary claim does not survive: under the
preliminary budget no operator improves on a constant predictor by more than 2.9\% on any shell of
1584 satellites or more; with a five-fold budget the heat kernel is better on the smallest shell on
both tasks; and at 1584 satellites the paired design resolves a difference 4.3 times the one under
audit. The preliminary direction is recovered only on the seam-cut graph, on nine seeds, on one of
two metrics, and a control that removes as many links at random positions reproduces it below its
own resolution, so its mechanism is not identified. Every claim is keyed to the table it is read
from, minimum detectable effects are printed beside every null, and which long runs collapse is
shown to depend on the floating-point precision.

\textbf{Why Computer Communications.} The contribution is an evaluation protocol and a negative
result about how propagation operators on satellite network graphs should be compared, in the line
of the journal's recent extended evaluations of protocols and reproducible experimental procedures,
rather than a new operator. The paper reports the regimes in which the audited method loses as
carefully as those in which it does not.

\textbf{Relation to prior work by the author.} The preliminary version was a case study inside a
manuscript that was reviewed and not accepted elsewhere; it is not published, and it is cited here
only as a preliminary version together with its released artifact, whose own headline this paper
shows cannot be reproduced from that artifact. No text, figure or result file is shared with any
manuscript of mine currently under review.

\textbf{Reproducibility.} All code, every result file of the 5590 runs, the scripts that produce
every figure, table and number, and the dated protocol declarations are released at
\url{https://github.com/haodpsut/leo-propagation-reach}, release tag \texttt{comcom-v1}.

The manuscript is single-authored, is not under consideration elsewhere, and declares the use of a
generative AI assistant for code and editing in a dedicated section, as Elsevier policy requires.
I declare no competing interests.

Sincerely,

Phuc Hao Do\\
CAIRA, Danang Architecture University, Da Nang, Viet Nam\\
\texttt{haodp@dau.edu.vn}
\end{document}
